\section{Euler inequalities for the deformation presentation}
\label{sec:deformation-euler}

We finish the two remaining inequalities in \eqref{eq:single-rees}.
The proof uses a positivity argument for graded Koszul complexes and a
short exact sequence which extracts the Smith exponents separately.

\subsection{Positivity for homogeneous sections}

Let \(D\) be a Noetherian local ring with maximal ideal \(\mathfrak m\).
Consider a standard \(\mathbb N^k\)-graded algebra \(R\) over \(D\),
generated by finitely many homogeneous elements in each unit-vector
degree \(e_i\). Let \(\mathfrak b=\prod_i(R_{e_i})R\), where the
product is of the corresponding generated ideals. For a finite graded
\(R\)-module \(P\), consider a list \(S\) of homogeneous sections
of degrees in \(\mathbb N^k\). Say the list satisfies the finiteness
condition on \(P\) if some power of \(\mathfrak m\mathfrak b\) is
contained in \(\operatorname{ann}P+(S)\).
For large multidegree \(d\), every finite graded module annihilated
by this latter ideal has finite \(D\)-length in degree \(d\).
Indeed monomials generating a given component, times one of finitely
many module generators, have at least \(n\) section factors of every
unit degree whenever all coordinates of \(d\) are sufficiently large.
Thus \(\mathfrak m^n\) kills that component if
\((\mathfrak m\mathfrak b)^n\) kills the module. The component is
finitely generated over \(D\), and \(D/\mathfrak m^n\) is Artinian.

Write \(\chi(S;P)_d\) for the alternating sum of the lengths, in degree
\(d\), of the homology of the graded Koszul complex on \(S\). Its
homology is finite as a module, and every section in \(S\) annihilates
it by the Koszul homotopy. The preceding observation therefore makes
the definition valid in all sufficiently large multidegrees. A short
exact sequence of coefficient modules yields additivity of \(\chi\)
by the homology sequence. A permutation of the list gives the signed
permutation isomorphism of Koszul complexes, preserving \(\chi\).
If \(s\) has degree \(p\), and \(S'\) is the rest of the list, then
\[
 \chi(s,S';P)_d=\chi(S';P/sP)_d
                     -\chi(S';P[s](-p))_d .
\]
Here the degree-\(d\) part of \(P(-p)\) is \(P_{d-p}\).
To see the equality, take first the homology of
\(P(-p)\xrightarrow{s}P\) in the Koszul double complex; the two rows
left are exactly the indicated kernel and cokernel with the
remaining sections. The Euler sums can be taken on the page after
their remaining Koszul homologies have been computed, where the
terms have finite length. This also connects the Koszul definition
with the recursive Euler characteristic in the source.

We prove \(\chi(S;P)_d\ge0\) for all sufficiently large \(d\).
Induct on the number of sections. The case of the empty list follows
from the definition of length. By additivity it suffices to handle
a shift of \(R/\mathfrak p\) for a homogeneous prime
\(\mathfrak p\supseteq\operatorname{ann}P\). Indeed there is a
finite homogeneous prime filtration of \(P\). Choose a nonzero
homogeneous element with maximal annihilator among such elements.
Its homogeneous annihilator is prime since, for homogeneous
\(a,b\), \(abx=0\) and \(bx\ne0\) imply that \(x\) and \(bx\) have
the same annihilator. Homogeneous primality is primality here by
using any translation-invariant total order on \(\mathbb Z^k\) and
the extremal nonzero homogeneous components of products in the
quotient. Repeat on the quotient by the generated cyclic module. The ascending chain condition terminates the filtration. The
finiteness condition on the original module holds for each factor
and is unaffected in large degrees by its grading shift.

If \(\mathfrak p\) contains \(\mathfrak b\), the large components of
the factor and of the finitely many shifts in its complex vanish.
Otherwise there is a section outside \(\mathfrak p\) in each unit
degree. Write \(P_0=R/\mathfrak p\) and take the first section \(s\).
If \(s\notin\mathfrak p\) it is regular on \(P_0\), so the recursion
gives \(\chi(S;P_0)_d=\chi(S';P_0/sP_0)_d\), nonnegative by
induction. If \(s\in\mathfrak p\), take \(z\notin\mathfrak p\) of
the same degree \(p\), by multiplying sections of the unit degrees
(the empty product is \(1\)). Now \(S'\) already satisfies the
finiteness condition on \(P_0\). The recursion for the zero action
of \(s\) gives
\(\chi(S';P_0)_d-\chi(S';P_0(-p))_d\). Additivity for the injection
\(P_0(-p)\xrightarrow{z}P_0\) identifies this difference with
\(\chi(S';P_0/zP_0)_d\), again nonnegative by induction.
The finite number of factors permits one common lower threshold for
the multidegree. This proves positivity. In particular passage to a
quotient can only decrease \(\chi(S;-)_d\) under the finiteness
condition on the original module, by additivity for its kernel.

The same recursion bounds the Euler characteristic by successively
discarding annihilator terms:
\[
      \chi(S,U;P)_d\le\chi(U;P/(S)P)_d
\]
when the combined list satisfies the finiteness condition on \(P\).
At each step \(P[s]\) and \(P/sP\) are killed by \(s\), so the
remaining list satisfies the condition on them. The subtracted
term is eventually nonnegative by the result just proved.

\subsection{Finiteness on the chart and the middle comparison}

Apply this with \(D=A_0\), \(R=\mathcal R\), \(S=s_*\), and the
trace lift \(a\) chosen in Section~\ref{sec:deformation-chart}.
Let \(\mathcal J\) be the coefficient ideal of series killed by
the chart and \(M=\mathcal R/\mathcal J\). The list \(S,a\)
satisfies the finiteness condition on \(M\). To see it, take a
prime containing \(\mathcal J+(S,a)\) which avoids \(\mathfrak b\).
It contains all \(Z_i\) because the chart kills them. By
\eqref{eq:rees-contract} its contraction to \(D\) contains \(I_D\).
A vertical contraction equals \(\mathfrak m_D\) by the finiteness
of \(D/I_D\). A horizontal contraction admits an embedding in
\(\overline F\). It cannot contain \(\Delta\) by
\eqref{eq:rees-boundary}. As proved in Section~\ref{sec:deformation-support},
this embedding extends to a mixed geometric solution of the
equations. Its value on \(a\) would be zero, contradicting the
nonvanishing assumption. Thus each prime in question contains
\(\mathfrak m_D\). Taking radicals and then a power gives the stated
finiteness condition.

The chosen \(\phi(a)\) is a nonzero series and \(a\) is regular on
\(M\). Indeed multiplication by \(a\) acts on each coefficient via
multiplication by \(\phi(a)\) in a domain. Taking the Euler sum first
in \(a\), then moving \(a\) to the end, and discarding the
annihilator terms for \(S\) gives
\[
 \chi(S;M/aM)_d=\chi(a,S;M)_d=\chi(S,a;M)_d
              \le\chi(a;M/(S)M)_d .
\]
This proves the middle inequality of \eqref{eq:single-rees}.

\subsection{Extracting the uniformizer powers}

Use the leading ideal \(\mathfrak q\) of
Section~\ref{sec:deformation-chart} and put \(G=\mathcal R/\mathfrak q\).
The kernel of the chart lies in the leading kernel in every degree.
Also \(a\in\mathfrak q\), since \(x_A(a)=0\) makes every monomial
in its degree-zero image have positive negative weight. Consequently
\(M/aM\) surjects onto \(G\). The quotient comparison gives
\[
                         \chi(S;G)_d\le\chi(S;M/aM)_d .
\]
The leading coefficient lattice in every degree is torsion-free,
so \(\varpi\) is regular on \(G\). The special fiber
\(\overline G=G/\varpi G\) has finite-length components, with the
explicit Hilbert function \(H(d)=\operatorname{length}_D\overline G_d\)
proved in Section~\ref{sec:deformation-chart}. For any list of
sections \(U\), the chain modules in its degree-\(d\) Koszul complex
on \(\overline G\) all have finite length. Taking the alternating
sum on the chains or on their homologies gives
\[
        \chi(U;\overline G)_d=
              \prod_{u\in U}(1-T_{\deg u})H(d),
       \qquad T_p H(d)=H(d-p).
\]

The classes of the first \(r\) sections of \(S\) on \(G\) are
\(\varpi^{2c_i} y_i^2 T_0\). In fact the difference
\(F_i-\varpi^{c_i}y_i\) has weight \(-1\) and belongs to the
saturated square, hence has negative filtration at least two.
The sum \(F_i+\varpi^{c_i}y_i\) has filtration at least one.
Their product has filtration at least three, so it vanishes on the
layer \(n=2\) of the first Rees component. This proves the claim.

Let \(y\) be any section of degree \(p\), \(c\ge0\), and let \(U\)
be a remaining list for which \(\varpi^c y,U\) satisfies the
finiteness condition on \(G\). The following extraction is exact.
\[
   \chi(\varpi^c y,U;G)_d
       =\chi(y,U;G)_d+c\,\chi(U;\overline G)_d .
\]
Indeed \(G[\varpi^c y]=G[y]\) by regularity, and there is a short
exact sequence
\[
   0\longrightarrow G/yG
       \xrightarrow{\ \varpi^c\ }G/(\varpi^c y)G
       \longrightarrow G/\varpi^cG\longrightarrow0 .
\]
Injectivity follows by cancelling \(\varpi^c\) in \(G\).
The final module has \(c\) successive factors isomorphic to
\(\overline G\), by the same cancellation. Additivity for \(U\)
and the Euler recursion prove the identity. Stripping a scalar
factor enlarges the ideal of sections, so the required finiteness
condition is preserved at subsequent steps.

Apply the identity in turn to the first \(r\) sections, moving each
into first position by Koszul permutation. In the term measured on
\(\overline G\), all the other \(r-1\) sections have the same degree
\(e_0\), regardless of the scalar powers already stripped. Thus each
extraction contributes \(2c_i\) times the same difference
\[
        \mathcal D H=(1-T_{e_0})^{r-1+s}
                        (1-T_{e_1})^{r'+s'}(1-T_{e_2})H.
\]
The Euler characteristic on \(G\) with all \(r\) scalar powers
stripped is eventually nonnegative. We conclude that
\[
 \chi(S;M/aM)_d\ge \chi(S;G)_d
                \ge\left(\sum_{i=1}^r2c_i\right)\mathcal D H(d).
\]
The generating-function computation in
Section~\ref{sec:single-character} gives
\(\mathcal DH(d)=2^{r-1+r'}\) for sufficiently large \(d\).
This proves the leftmost inequality in \eqref{eq:single-rees}.
Combining it with the middle comparison above and the bound of
Section~\ref{sec:deformation-support} proves the full graded
intersection estimate for the stated model.

\subsection{The cut chart and finiteness of the set of orbits}

We give the last two geometric constructions used above from the
model hypotheses. Before cutting, smoothness in the source gives
a graded chart \(\phi_0\) over \(F\) with \(m\) variables of weight
zero, in addition to the \(r,s,r'+1,s'+1\) variables of nonzero
weights already specified. Its linear terms separate the axes, and
an augmentation element whose image has zero linear term belongs
to the saturated square. Thus the map from the saturated cotangent
lattice of weight \(-1\) to the vector space of linear coefficients
on the \(r\) weight-\(-1\) variables is injective. Every coordinate
axis is met by the image, up to a nonzero scalar, by projecting an
axis-separating element to weight \(-1\). This proves that the
lattice has rank \(r\). Its finite generation over \(\mathcal O\)
follows from finite generation of \(A_{-1}\) over \(D\): on the
cotangent classes \(D\) acts through the augmentation, since
\((b-x_A(b))h\in(\ker x_A)^2\) for every \(h\in A_{-1}, b\in D\).
The weight-zero analogue, applied after subtracting constant terms,
shows that the linear coefficient vectors of the \(m\) equations
\(Z_i\) span all \(m\) zero-weight directions because those equations
span their saturated cotangent lattice over \(F\). They therefore
give an invertible square matrix of linear terms in the zero-weight
variables. There are no linear terms of other weights.

In the formal series ring over \(F\), replace the \(m\) zero-weight
variables by \(\phi_0(Z_i)\), keeping the nonzero-weight variables.
This substitution has zero constant part and invertible linear part.
It has a unique inverse by formal degree recursion: solve the
linear part first, and at each subsequent degree apply its inverse
matrix to the error in that homogeneous degree. Each coefficient
stabilizes at its own degree. The substitution preserves weights, and
so does its inverse by the same recursion. Compose the original
chart with this inverse and then set all zero-weight variables to
zero. This kills the \(Z_i\), leaves the constant coefficient equal
to \(x_A\), and preserves the nonzero-weight linear terms, because
the linear substitution was block diagonal and the nonzero block
was the identity. This is exactly the cut chart used in the leading
coefficient calculation.

For orbit finiteness take finite \(D\)-module generating sets of
\(A_{-2}\) and \(A_2\), whose existence was shown in
Section~\ref{sec:deformation-chart}. Every mixed point lies where
some product \(fg\) of such generators is nonzero. On this locus
each orbit has a point with \(f=1\), since under torus scaling \(f\) is
multiplied by \(z^{-2}\), and the base field is algebraically closed.
The slice algebra on \(f=1\), \(g\ne0\), after base change to
\(\overline F\), was shown there to have a field as every local ring
at a maximal ideal. It is Noetherian of dimension zero. To see the
dimension assertion, any prime is contained in a maximal ideal and
the localization there is a field, so the two primes coincide.
A zero-dimensional Noetherian ring is Artinian and has finitely many
maximal ideals. Thus this slice has finitely many geometric points.
The finitely many pairs \((f,g)\) cover the mixed locus, giving
finitely many mixed orbits and the desired finite set of
representatives \(T\).

The source comparisons occur in \(\mathrm{CutHolds}\),
\(\mathrm{LaDifPerm}\), \(\mathrm{AdmissibleMonoIdeal}\), and
\(\mathrm{AlgMemCoef}\). This completes the graded intersection
calculation from the smooth model, the cotangent spanning conditions,
finiteness over the coefficient DVR in weight zero, and the two
geometric point conditions. Constructing the arithmetic model and
establishing those conditions for its Galois points use the
deformation and automorphy theorems specified in
Section~\ref{sec:single-character}.
